\documentclass[twocolumn, 10pt]{article}


% --- Page Layout --- %

\usepackage[utf8]{inputenc}
\usepackage{xcolor}
\usepackage{geometry}
\geometry{
	a4paper,
	left=20.0mm,
	right=20.0mm,
	top=20.0mm,
	bottom=22.0mm,
	columnsep=5.0mm,
	headsep=6.5mm,
	footskip=12.5mm
}
\usepackage[aux]{rerunfilecheck}
\usepackage{float}
\usepackage{dblfloatfix}


% --- Modern Font Setup --- %

\usepackage{fontspec}
\usepackage{unicode-math}
\setmainfont{TeX Gyre Termes}
\setmathfont{TeX Gyre Termes Math}


% --- Math Symbols --- %

\usepackage{amsmath}
\usepackage{mathtools}


% --- Units --- %

\usepackage[
	locale=US,
	per-mode = power,
	exponent-mode = input,
	print-unity-mantissa = false,
	range-units = single
]{siunitx}

\DeclareSIUnit{\year}{yr}
\DeclareSIUnit{\erg}{erg}
\DeclareSIUnit{\parsec}{pc}
\DeclareSIUnit{\angstrom}{\AA}
\DeclareSIUnit{\solarmass}{\ensuremath{M_\odot}}

\DeclareSIUnit{\dex}{dex}
\DeclareSIUnit{\rms}{rms}
\DeclareSIUnit{\cgs}{cgs}


% --- Aesthetic Improvements --- %

\usepackage{graphicx}
\usepackage{booktabs}
\usepackage{microtype}
\usepackage[autostyle]{csquotes}


% --- Draft Figures --- %

\usepackage[
	filename=false, 
	allfiguresdraft=draft
]{draftfigure}


% --- Float Improvements --- %

\usepackage{float}
\floatplacement{figure}{htbp}
\floatplacement{table}{htbp}

\usepackage[
	section,
	below,
]{placeins}


% --- Abstract Adjustments --- %

\usepackage{abstract}
\renewcommand{\abstracttextfont}{\small}
\setlength{\absleftindent}{0pt}
\setlength{\absrightindent}{0pt}
\setlength{\abstitleskip}{0.25\baselineskip}
\AfterEndEnvironment{abstract}{\vspace{1.5\baselineskip}}


% --- Section Styling --- %

\usepackage{titlesec}
\titleformat{\section}{\normalfont\large\bfseries}{\thesection}{1ex}{}
\titleformat{\subsection}{\normalfont\normalsize\bfseries}{\thesubsection}{1ex}{}


% --- Caption Styling --- %
\usepackage[
	labelfont=bf,
	font=small,
]{caption}
\usepackage{subcaption}

\captionsetup[figure]{width=0.9\columnwidth}
\captionsetup[figure*]{width=0.9\textwidth}

% --- Header & Footer --- %

\usepackage{fancyhdr}
\pagestyle{fancy}

\fancyhf{}
\renewcommand{\headrulewidth}{0pt}
\fancyfoot[C]{\itshape\thepage}

\fancypagestyle{plain}{
  \fancyhf{}
  \fancyfoot[C]{\itshape\thepage}
  \renewcommand{\headrulewidth}{0pt}
}


% --- Citation Setup --- %

\usepackage[american]{babel}

\usepackage[
	colorlinks=true, 
	linkcolor=black,      
	citecolor=black,      
	urlcolor=black
]{hyperref}

\usepackage[noabbrev, nameinlink]{cleveref}
\newcommand{\bcref}[1]{\textbf{\cref{#1}}}
\newcommand{\bCref}[1]{\textbf{\Cref{#1}}}

\usepackage[
	backend=biber,
	style=authoryear,
	giveninits=true,   
	maxnames=2,        
	minnames=1,
	uniquename=false,
	uniquelist=false
]{biblatex}

\DeclareNameAlias{author}{family-given} 
\DeclareDelimFormat{multinamedelim}{\addspace\&\addspace}
\DeclareDelimFormat{finalnamedelim}{\addspace\&\addspace}
\DefineBibliographyStrings{english}{andothers={et al.}} 

\DeclareFieldFormat{labeldate:extrayear}{}

\DeclareCiteCommand{\cite}[\mkbibparens]
	{\usebibmacro{prenote}}%
	{\usebibmacro{citeindex}%
	 \printtext[bibhyperref]{%
	 	\printnames{labelname}%
	 	\setunit{\addspace}%
	 	\printfield{labelyear}%
	 	\printfield{extrayear}%
	 }%
	}%
	{\multicitedelim}%
	{\usebibmacro{postnote}}%

\DeclareDelimFormat{multicitedelim}{\addcomma\addspace}

\DeclareFieldFormat[article,periodical]{journaltitle}{\textit{#1}}

\DeclareFieldFormat[article]{title}{}
\DeclareFieldFormat[book]{title}{}

\AtEveryBibitem{%
	\clearfield{volume}%
	\clearfield{number}%
	\clearfield{pages}%
	\clearfield{note}%
	\clearfield{location}%
	\clearfield{address}%
}

\renewcommand{\finentrypunct}{}

\DeclareBibliographyDriver{article}{%
	\usebibmacro{bibindex}%
	\usebibmacro{begentry}%
	\iffieldundef{doi}
		{\iffieldundef{url}
			{\printnames{author}\setunit{\addspace}%
			 \usebibmacro{journal}\setunit{\addspace}%
			 \printtext{\mkbibparens{\printfield{labelyear}\printfield{extrayear}}}}%
			{\href{\thefield{url}}{%
				\printnames{author}\setunit{\addspace}%
				\usebibmacro{journal}\setunit{\addspace}%
				\printtext{\mkbibparens{\printfield{labelyear}\printfield{extrayear}}}}}}
		{\href{https://doi.org/\thefield{doi}}{%
			\printnames{author}\setunit{\addspace}%
			\usebibmacro{journal}\setunit{\addspace}%
			\printtext{\mkbibparens{\printfield{labelyear}\printfield{extrayear}}}}}%
	\usebibmacro{finentry}%
}

\addbibresource{source.bib}


% --- Paragraph Skip --- %

\usepackage{parskip}
\setlength{\parindent}{0pt}
\setlength{\parskip}{0.4\baselineskip plus 2pt minus 2pt}


% --- Linted Code --- %

\usepackage{varwidth}
\usepackage{listings}
\usepackage{accsupp}

\newcommand{\noncopynumber}[1]{%
     \BeginAccSupp{ActualText={}}#1\EndAccSupp{}%
}

\definecolor{firebrick}{RGB}{178, 34, 34}     
\definecolor{goldenrod}{RGB}{218, 165, 32}    
\definecolor{olivedrab}{RGB}{107, 142, 35}    
\definecolor{steelblue}{RGB}{70, 130, 180}    
\definecolor{rebeccapurple}{RGB}{102, 51, 153}

\definecolor{lightergray}{RGB}{250, 250, 250} 
\definecolor{lightishgray}{RGB}{180, 180, 180}

\lstset{
	language=SQL,                          
	backgroundcolor=\color{lightergray},   
	basicstyle=\ttfamily\footnotesize\color{black},
	identifierstyle=\color{black},
	keywordstyle=\color{steelblue}\bfseries, 
	stringstyle=\color{goldenrod},        
	morekeywords=[2]{TOP, APPLY},
	keywordstyle=[2]{\color{firebrick}\bfseries},
	numberstyle=\tiny\color{lightishgray}\noncopynumber, 
	numbers=left,                          
	numbersep=8pt,                         
	breaklines=true,                       
	breakatwhitespace=false,               
	showstringspaces=false,                
	keepspaces=true,
	frame=single,                          
	rulecolor=\color{black},
	tabsize=4,                             
	keywordsprefix=,
	literate=*
		{=}{{{\color{rebeccapurple}=}}}1
		{!=}{{{\color{rebeccapurple}!=}}}2
		{<}{{{\color{rebeccapurple}<}}}1
		{>}{{{\color{rebeccapurple}>}}}1
		{+}{{{\color{rebeccapurple}+}}}1
		{-}{{{\color{rebeccapurple}-}}}1
		{*}{{{\color{rebeccapurple}*}}}1
		{:}{{{\color{rebeccapurple}:}}}1
		{10000}{{{\color{olivedrab}10000}}}5
		{\ 1\ }{{{\ \color{olivedrab}1\ }}}3
		{\ 3}{{{\ \color{olivedrab}3}}}2
		{0\.35}{{{\color{olivedrab}0.35}}}4
		{\ 0\.0}{{{\ \color{olivedrab}0.0}}}4
		{5\.0}{{{\color{olivedrab}5.0}}}3
		{2\.355}{{{\color{olivedrab}2.355}}}5
		{\ 500}{{{\ \color{olivedrab}500}}}4
}


% --- Dummy Text --- %

\usepackage{lipsum}


% --- Writing Lengths --- %

\AtEndDocument{
	\directlua{
		local factor = 65536 * 72.27
		
		local tw = math.floor((tex.getdimen("textwidth") / factor) * 100 + 0.5) / 100
		local cw = math.floor((tex.getdimen("columnwidth") / factor) * 100 + 0.5) / 100
		local cs = math.floor((tex.getdimen("columnsep") / factor) * 100 + 0.5) / 100
		local lw = math.floor((tex.getdimen("linewidth") / factor) * 100 + 0.5) / 100
		
		texio.write_nl("===================================================")
		texio.write_nl("CURRENT LAYOUT:")
		texio.write_nl("Textwidth:  " .. tw .. " in")
		texio.write_nl("Columnwidth: " .. cw .. " in")
		texio.write_nl("Columnsep:   " .. cs .. " in")
		texio.write_nl("Linewidth:   " .. lw .. " in")
		texio.write_nl("===================================================")
		io.flush()
	}
}

\setmainfont{Latin Modern Roman}
\setmathfont{Latin Modern Math}

\title{\Large\textbf{Revisiting Neutrinos via Charm Decays in Astrophysical Sources}}%
\author{Fritz Ali Agildere%
		\vspace{-0.5ex}\\\small\href{mailto:fritz.agildere@udo.edu}{fritz.agildere@udo.edu}%
		\\\small\textit{Erasmus Mundus Joint Master in Astrophysics and Space Science}}%
\date{\small \today}

\fancyhead[C]{\textit{Revisiting Neutrinos via Charm Decays in Astrophysical Sources}}
\fancyfoot[R]{\textit{Master in Astrophysics and Space Science}}
\fancyfoot[L]{\textit{Bachelor Thesis Addendum and Revisions}}
\fancypagestyle{plain}{
	\fancyhead[C]{}
	\fancyfoot[R]{\textit{Master in Astrophysics and Space Science}}
	\fancyfoot[L]{\textit{Bachelor Thesis Addendum and Revisions}}
}

\hypersetup{
	pdfauthor={Fritz Ali Agildere},
	pdftitle={Revisiting Neutrinos via Charm Decays in Astrophysical Sources}
}

\ExecuteBibliographyOptions{maxbibnames=1, minbibnames=1}
\graphicspath{{../plots/}}
\usepackage{multicol}

\DeclareNameFormat{first}{\ifnumequal{\value{listcount}}{1}{\namepartfamily}{}}
\DeclareCiteCommand{\cite}[\mkbibparens]
	{\usebibmacro{prenote}}%
	{\usebibmacro{citeindex}%
	 \printtext[bibhyperref]{\printnames[first][1-1]{labelname}\setunit{\addspace}\printfield{labelyear}\printfield{extrayear}}}%
	{\multicitedelim}%
	{\usebibmacro{postnote}}%
\AtEveryCitekey{\renewbibmacro*{name:andothers}{}}

\newcommand{\fig}[4]{%
	{\setlength{\intextsep}{3pt plus 1pt minus 1pt}\captionsetup[figure]{skip=2pt, width=#4\textwidth}%
	\begin{figure}[H]
		\centering
		\includegraphics[width=#4\textwidth, draft=false]{#1}
		\caption{#2}
		\label{#3}
	\end{figure}}}

\newcommand{\figsection}[1]{\vspace{0.25\baselineskip}{\centering\normalfont\large\bfseries #1\par}\vspace{0.25\baselineskip}}

\raggedbottom
\hyphenpenalty=10000
\exhyphenpenalty=10000
\emergencystretch=2em

\begin{document}

	\onecolumn
		\maketitle
		\begin{abstract}
			A recheck of the code behind my bachelor thesis found a unit error in the conversion of cross sections, a mismatched normalization
			of the charm cross section and an unbounded optical depth, and a vectorized rebuild with improved integration runs about sixty
			times faster. With these corrected, magnetar results hold and charm separates more clearly from pions and kaons, while charm
			in the accretion disk takes over at much higher energies than stated in the thesis.
		\end{abstract}

	\begin{multicols}{2}

	\section{Corrections}

	The thesis code was rebuilt as one vectorized program. It reproduces charm spectra to \num{1e-6} relative precision using the thesis settings,
	while every parametrization agrees with its source and the method itself stands.

	\textbf{Units:} Cooling factor and optical depth convert cross sections from millibarn to square centimetres with the factor that belongs to
	barn. Each target therefore acted a thousand times denser than stated.

	\textbf{Charm:} In the thesis code, the charm cross section was divided by an extra factor of roughly \num{22} after its integral over the
	entire energy fraction range was compared to data from forward measurements \cite{aaij2016}. Most of that integral however comes from energy
	fractions far below the measured window, and inside the window the parametrization \cite{gonc2007} already gives \num{1.07} times the data,
	so no extra factor was needed.

	\textbf{Depth:} The optical depth multiplied spectra without bound, although the share of interacting protons cannot exceed unity.
	The corrected model uses the interaction probability instead, matching the thesis variant without optical depth as long as
	the ejecta stay opaque.

	\textbf{Numerics:} Trapezoid weights replace the first order sum, a stable exponential removes a cancellation in the cooling factor and
	the vectorized integration runs about sixty times faster.

	\section{Magnetar}

	All figures compare the thesis on the left with the corrected results on the right and are normalized to the charm peak, shaded bands mark
	a factor of three around charm, and the reference \cite{carp2020} given in \bCref{fig:ref} shows pions and kaons with and without muon decay.
	In the light curves depicted in \bCref{fig:light} charm still peaks first followed by kaons and pions, but all cooling features move to earlier
	times by about one decade. Charm now peaks after \qty{7.1e3}{\second} shortly after the spindown time of \qty{3.2e3}{\second} as is described in
	the reference, instead of after \qty{4.9e4}{\second} in the thesis.

	The charm composition in \bCref{fig:charm} holds, with $D^0$ and $D^+$ carrying shares of \num{0.51} and \num{0.38} at the maximum and
	the individual peaks ordered by lifetime.

	In the fluence of \bCref{fig:fluence} charm exceeds both meson components above \qty{2.9e9}{\GeV} and leads the larger one at
	\qty{1e10}{\GeV} by a factor of \num{14.5} instead of \num{3.0} so that even the lower edge of its band clears pions and kaons. The
	discrepancy that the thesis attributed to the optical depth disappears with the unit correction, and the initial proton energy of the
	thesis remains the correct evaluation of the formula given in the reference.

	\section{Accretion Disk}

	The disk result changes. Still, the charm composition in \bCref{fig:diskcharm} keeps its shape. However, at the stated density of
	\qty{1e14}{\per\cubic\centi\metre} pions instead of kaons dominate the low energies in \bCref{fig:diskfluence} and charm
	takes over only above \qty{1.6e10}{\GeV} instead of \qty{3.7e8}{\GeV} in the thesis. At \qty{1e11}{\GeV} it leads by a
	factor of \num{2.3} instead of \num{6.8} so that its lower band edge stays below the kaons. As \bCref{fig:density} shows,
	the thesis takeover energy returns only near \qty{4.6e15}{\per\cubic\centi\metre} which lies more than one decade above
	the stated density.

	\section{Conclusion}

	Overall, the thesis method stands while the magnetar conclusions become stronger, with the disk
	conclusions holding only for a much denser target.

	\vspace{1ex}
	\printbibliography[heading=none]

	\end{multicols}

	\clearpage

	\figsection{Magnetar Scenario}

	\fig{referenceMagnetar.pdf}{Reference light curve at $E_\nu = \qty{1e9}{\GeV}$ and fluence \cite{carp2020} normalized to the charm peak.}{fig:ref}{0.942}

	\fig{comparedMagnetarLightCurve.pdf}{Light curves at $E_\nu = \qty{1e9}{\GeV}$ normalized to the charm maximum.}{fig:light}{0.942}

	\fig{comparedMagnetarFluence.pdf}{Fluence normalized to the charm peak of the widest time window.}{fig:fluence}{0.942}

	\fig{comparedMagnetarCharm.pdf}{Charm hadrons in the light curve normalized to the total charm maximum.}{fig:charm}{0.942}

	\figsection{Accretion Disk Scenario}

	\fig{comparedDiskCharm.pdf}{Charm hadrons in the disk spectrum normalized to the total charm maximum.}{fig:diskcharm}{0.942}

	\fig{comparedDiskFluence.pdf}{Disk spectra by parent normalized to the charm maximum.}{fig:diskfluence}{0.942}

	\fig{correctedDiskDensity.pdf}{Neutrino energy at which charm overtakes pions and kaons in the disk against density after the corrections,
	compared with the thesis value at its stated density.}{fig:density}{0.538}

	\vfill

	\noindent\small\textit{Declaration:} The author acknowledges the use of artificial intelligence products by Anthropic, specifically of the
	models Claude 5.5 Sonnet and Claude 5.5 Opus for rechecking and rebuilding the code, preparing the figures and drafting this summary. The
	author independently verified all outcomes and maintains full accountability for this document.

\end{document}
